\section{Hard sync 的精确几何与稳定区间}
\label{sec:hard-sync-geometry}

本节固定一个 prefix，省略条件变量。沿用 $\pi_k$ 表示第 $k$ 轮冻结 anchor，$\pi_T$ 表示固定 teacher。词表 $\mathcal V$ 有限，$\pi_0(v)>0$、$\pi_T(v)>0$ 对所有 token 成立。假设 inner 优化确实得到 Gibbs target，随后将这个 target 完整复制为下一轮 anchor。这个假设同时排除了有限 inner 误差、模型表达能力不足和参数同步与分布同步之间的差异。

\begin{definition}[精确 Gibbs hard sync]
对任意 $\eta>0$ 和严格正分布 $\pi$，定义
\[
 T_\eta(\pi)(v)
 =\frac{\pi(v)^{1-\eta}\pi_T(v)^\eta}{Z_\eta(\pi)},
 \qquad
 Z_\eta(\pi)=\sum_{w\in\mathcal V}\pi(w)^{1-\eta}\pi_T(w)^\eta.
\]
精确 hard sync 是 $\pi_{k+1}=T_{\eta_k}(\pi_k)$。$\eta_k$ 是 mirror 更新的步长；它与优化器学习率及每轮 inner steps 是不同量。
\end{definition}

严格正性使负指数也有定义。因此本节先分析所有 $\eta>0$ 的数学更新，再区分通常采用的 $0<\eta\leq1$ 区间。若初始分布有零概率，$0<\eta<1$ 的更新会一直保留这些零概率，不能恢复 teacher 的完整支持集；$\eta>1$ 的负幂则在这些 token 上失去定义。有限 logits 的 softmax 满足正性，而显式截断后的采样分布不能直接代入本节结论。

\subsection{归一化常数消去后的精确收缩}

\begin{definition}[Log-ratio 跨度与 Hilbert 距离]
对 token 函数 $f$，定义
\[
 \operatorname{osc}_v f_v=\max_v f_v-\min_v f_v.
\]
对两个严格正概率分布，定义
\[
 d_H(\pi,\widetilde\pi)
 =\operatorname{osc}_v\log\frac{\pi(v)}{\widetilde\pi(v)},
 \qquad b_k=d_H(\pi_k,\pi_T).
\]
$d_H$ 是正概率单纯形上的 Hilbert 距离。它忽略 log-density 中与 token 无关的加性常数，恰好对应概率归一化所消除的自由度。
\end{definition}

这个定义确实给出距离。非负性来自最大值减最小值；对称性来自 $\operatorname{osc}(-f)=\operatorname{osc}(f)$；三角不等式来自 $\operatorname{osc}(f+g)\leq\operatorname{osc}(f)+\operatorname{osc}(g)$。若 $d_H(\pi,\widetilde\pi)=0$，则两个分布的概率比是常数，归一化迫使这个常数等于一。

\begin{theorem}[两条 hard-sync 轨迹的精确收缩]
\label{thm:hs-pairwise-contraction}
相同 teacher 和相同步长下，任意严格正分布满足
\begin{equation}\label{eq:hs-pairwise-contraction}
 d_H(T_\eta(\pi),T_\eta(\widetilde\pi))
 =|1-\eta|d_H(\pi,\widetilde\pi).
\end{equation}
特别地，$\pi_T$ 是不动点，而且
\[
 b_{k+1}=|1-\eta_k|b_k.
\]
\end{theorem}

\begin{proof}
先逐 token 相除，再取对数：
\[
 \log\frac{T_\eta(\pi)(v)}{T_\eta(\widetilde\pi)(v)}
 =(1-\eta)\log\frac{\pi(v)}{\widetilde\pi(v)}
 -\log Z_\eta(\pi)+\log Z_\eta(\widetilde\pi).
\]
最后两个项与 token 无关，在跨度中抵消。对任意实数 $c$，$\operatorname{osc}(cf)=|c|\operatorname{osc}(f)$，从而得到式~\eqref{eq:hs-pairwise-contraction}。将 $\widetilde\pi=\pi_T$ 代入，并使用 $T_\eta(\pi_T)=\pi_T$，得到第二个结论。
\end{proof}

\begin{theorem}[常步长与变步长的精确解]
\label{thm:hard-sync}
记
\[
 A_0=1,\qquad A_k=\prod_{j=0}^{k-1}(1-\eta_j),
 \qquad u_v=\log\frac{\pi_0(v)}{\pi_T(v)}.
\]
任意正步长序列的精确解为
\begin{equation}\label{eq:hard-closed}
 \pi_k(v)=\frac{\pi_T(v)e^{A_k u_v}}
 {\sum_w\pi_T(w)e^{A_k u_w}},
 \qquad b_k=|A_k|b_0.
\end{equation}
常步长时 $A_k=(1-\eta)^k$，也可以写成
\[
 \pi_k(v)=\frac{\pi_0(v)^{(1-\eta)^k}
 \pi_T(v)^{1-(1-\eta)^k}}
 {\sum_w\pi_0(w)^{(1-\eta)^k}
 \pi_T(w)^{1-(1-\eta)^k}}.
\]
\end{theorem}

\begin{proof}
$k=0$ 时 $\pi_T(v)e^{u_v}=\pi_0(v)$，配分函数等于一。假设第 $k$ 轮公式成立，其归一化常数记为 $C_k$。下一轮未归一化权重为
\begin{align*}
 \pi_k(v)^{1-\eta_k}\pi_T(v)^{\eta_k}
 &=C_k^{-(1-\eta_k)}\pi_T(v)
       e^{(1-\eta_k)A_k u_v}\\
 &=C_k^{-(1-\eta_k)}\pi_T(v)e^{A_{k+1}u_v}.
\end{align*}
第一项不依赖 token，归一化后消失。由归纳得到分布闭式。取相对 teacher 的 log-ratio 后，归一化常数同样在跨度中抵消，得到 $b_k=|A_k|b_0$。
\end{proof}

\begin{corollary}[常步长的完整稳定分类]
对非 teacher 初值，精确 hard sync 的行为如下。
\begin{enumerate}
 \item $0<\eta<1$ 时，$A_k>0$ 且趋于零，策略从同一侧趋向 teacher。
 \item $\eta=1$ 时，第一轮更新等于 teacher，此后保持不变。
 \item $1<\eta<2$ 时，$A_k$ 正负交替并趋于零，策略在 teacher 两侧交替收敛。
 \item $\eta=2$ 时，$A_k=(-1)^k$，形成精确二周期。
 \item $\eta>2$ 时，$|A_k|$ 发散，log-ratio 跨度发散，策略不能趋向 teacher。
\end{enumerate}
因此完整的全局渐近稳定区间是 $0<\eta<2$。teacher KL 的每步单调下降需要另外分析。
\end{corollary}

\begin{proof}
前四项直接由式~\eqref{eq:hard-closed} 得到。第五项中，偶数轮的 $A_k$ 趋于正无穷，奇数轮趋于负无穷。偶数轮的概率质量集中到 $u_v$ 最大的 token 集合，奇数轮集中到最小的 token 集合；各集合内部的极限概率与 $\pi_T$ 成比例。非 teacher 初值使 $u$ 非常数，所以两个集合互不相交。

有限正 teacher 下仍有
\[
 \KL(\pi_k\Vert\pi_T)\leq\log\frac1{\min_v\pi_T(v)}.
\]
因此这里发散的是 log-ratio 跨度和反向 KL，而正向 teacher KL 保持有界。反向 KL 发散，是因为上述集中使某些具有正 teacher 质量的 token 概率趋于零。
\end{proof}

\begin{proposition}[变步长的充要条件]
非 teacher 初值下，$\pi_k\to\pi_T$ 当且仅当 $A_k\to0$。若每个 $\eta_k\in(0,2)$，且没有任何一步等于一，则这个条件等价于
\[
 \sum_{k=0}^{\infty}\eta_k(2-\eta_k)=\infty.
\]
若某一步等于一，则此后 $A_k=0$，直接收敛。
\end{proposition}

\begin{proof}
$A_k\to0$ 时由有限词表闭式得到逐 token 收敛。反过来，概率趋向严格正 teacher 时 $b_k\to0$；因为 $b_0>0$，故 $|A_k|\to0$。对 $a_k=|1-\eta_k|\in(0,1)$，无限乘积 $\prod_k a_k$ 为零，当且仅当 $\sum_k(1-a_k)$ 发散。这个标准乘积判据可以由 $-\log a\geq1-a$，以及 $a$ 接近一时 $-\log a\leq2(1-a)$ 推出。最后使用
\[
 1-a_k\leq1-a_k^2\leq2(1-a_k),
 \qquad 1-a_k^2=\eta_k(2-\eta_k),
\]
得到等价条件。
\end{proof}

例如 $\eta_k=1/(k+2)$ 给出 $A_K=1/(K+1)\to0$。若改成 $\eta_k=1/(k+2)^2$，则
\[
 A_K=\prod_{n=2}^{K+1}\frac{(n-1)(n+1)}{n^2}
 =\frac{K+2}{2(K+1)}\longrightarrow\frac12.
\]
第二种步长虽然每轮都执行正确方向的 hard sync，但总更新量有限，最终仍保留初始 log-ratio 的一半。

\subsection{配分函数与两个方向的 KL 路径}

沿用上节 $u_v$，定义一条由初始策略和 teacher 唯一确定的路径：
\[
 \psi(a)=\log\sum_v\pi_T(v)e^{a u_v},\qquad
 \pi^{[a]}(v)=\pi_T(v)e^{a u_v-\psi(a)}.
\]
这里 $a$ 是路径坐标，实际第 $k$ 轮取 $a=A_k$；它不代表新的策略变量。归一化给出 $\psi(0)=\psi(1)=0$，以及 $\pi^{[0]}=\pi_T$、$\pi^{[1]}=\pi_0$。

\begin{theorem}[KL 路径的精确积分]
\label{thm:hs-kl-path}
配分函数满足
\[
 \psi'(a)=\E_{\pi^{[a]}}u,
 \qquad\psi''(a)=\Var_{\pi^{[a]}}(u).
\]
沿路径的正向与反向 KL 分别为
\begin{align}
 \KL(\pi^{[a]}\Vert\pi_T)
 &=a\psi'(a)-\psi(a)
 =\int_0^a z\psi''(z)\,\mathrm dz,
 \label{eq:hs-forward-kl-path}\\
 \KL(\pi_T\Vert\pi^{[a]})
 &=\psi(a)-a\psi'(0)
 =\int_0^a(a-z)\psi''(z)\,\mathrm dz.
 \label{eq:hs-reverse-kl-path}
\end{align}
任意路径两点之间还有
\[
 \KL(\pi^{[a]}\Vert\pi^{[c]})
 =(a-c)\psi'(a)-\psi(a)+\psi(c).
\]
\end{theorem}

\begin{proof}
有限词表允许逐项微分。令 $C(a)=\sum_v\pi_T(v)e^{au_v}$，则
\[
 \psi'(a)=\frac{C'(a)}{C(a)}=\sum_v\pi^{[a]}(v)u_v,
 \qquad
 \psi''(a)=\frac{C''(a)}{C(a)}-
          \left(\frac{C'(a)}{C(a)}\right)^2.
\]
右侧正好是 $u$ 的方差。正向 log-ratio 是 $au_v-\psi(a)$，在 $\pi^{[a]}$ 下取期望得到第一式。其关于 $a$ 的导数为 $a\psi''(a)$，在零点的值为零，积分得到式~\eqref{eq:hs-forward-kl-path}。反向 log-ratio 是 $\psi(a)-au_v$，在 teacher 下取期望得到第二式；再对 $\psi$ 在零点使用带积分余项的二阶展开，得到式~\eqref{eq:hs-reverse-kl-path}。最后两条路径分布的 log-ratio 为 $(a-c)u_v-\psi(a)+\psi(c)$，取期望即得。
\end{proof}

由于 $\psi''(a)\geq0$，正向 KL 在 $a>0$ 上随 $a$ 增大而增大，在 $a<0$ 上随 $a$ 增大而减小。反向 KL 的导数 $\psi'(a)-\psi'(0)$ 也具有相同符号。因此 $0<\eta\leq1$ 下，两个方向的 KL 都逐轮下降。$1<\eta<2$ 时，连续两轮的 $a$ 符号不同，跨度收缩本身不能比较它们的 KL 大小。

\begin{corollary}[局部 KL 速率与全局绝对界]
非 teacher 初值、固定 $0<\eta<2$ 且 $\eta\neq1$ 时，
\[
 \lim_{k\to\infty}
 \frac{\KL(\pi_{k+1}\Vert\pi_T)}{\KL(\pi_k\Vert\pi_T)}
 =(1-\eta)^2.
\]
对所有实数 $A_k$，还有
\[
 \KL(\pi_k\Vert\pi_T)\leq\frac{A_k^2b_0^2}{8}.
\]
\end{corollary}

\begin{proof}
在零点展开得到
\[
 \KL(\pi^{[a]}\Vert\pi_T)
 =\frac{a^2}{2}\Var_{\pi_T}(u)+O(|a|^3).
\]
非 teacher 初值使该方差严格正；将 $a=A_k$ 与 $A_{k+1}=(1-\eta)A_k$ 代入并相除，得到极限。任意支撑在长度为 $b_0$ 的区间内的随机变量，方差不超过 $b_0^2/4$。把这个界代入式~\eqref{eq:hs-forward-kl-path}，得到 $a^2b_0^2/8$；当 $a<0$ 时反向积分同样给出正的 $a^2/2$ 因子。
\end{proof}

这个全局绝对界以初始跨度衡量误差；它没有把当前 KL 与初始 KL 作比例比较。局部极限 $(1-\eta)^2$ 也不能提升为全局单步收缩因子。

\subsection{精确三点恒等式与单步下降}

\begin{theorem}[Hard sync 的三点 KL 恒等式]
\label{thm:hs-three-point}
记 $\pi^+=T_\eta(\pi)$。对任意 $\eta>0$，
\begin{equation}\label{eq:hs-three-point}
 \eta\KL(\pi\Vert\pi_T)
 =\eta\KL(\pi^+\Vert\pi_T)
 +(1-\eta)\KL(\pi^+\Vert\pi)
 +\KL(\pi\Vert\pi^+).
\end{equation}
\end{theorem}

\begin{proof}
从更新定义分别减去 anchor 和 teacher 的 log-density：
\begin{align*}
 \log\frac{\pi^+}{\pi}
 &=-\eta\log\frac{\pi}{\pi_T}-\log Z_\eta(\pi),\\
 \log\frac{\pi^+}{\pi_T}
 &=(1-\eta)\log\frac{\pi}{\pi_T}-\log Z_\eta(\pi).
\end{align*}
第一式在 $\pi$ 下取负期望，得到
\[
 \KL(\pi\Vert\pi^+)
 =\eta\KL(\pi\Vert\pi_T)+\log Z_\eta(\pi).
\]
第二式在 $\pi^+$ 下取期望，并与第一式在 $\pi^+$ 下的期望组合，得到
\[
 \eta\KL(\pi^+\Vert\pi_T)
 +(1-\eta)\KL(\pi^+\Vert\pi)
 =-\log Z_\eta(\pi).
\]
两式相加，配分函数消去，即为所需恒等式。
\end{proof}

当 $0<\eta\leq1$ 时，式~\eqref{eq:hs-three-point} 的最后两个项非负。常步长的下降量和累计更新量因此满足
\begin{align*}
 \KL(\pi_k\Vert\pi_T)-\KL(\pi_{k+1}\Vert\pi_T)
 &=\frac{1-\eta}{\eta}\KL(\pi_{k+1}\Vert\pi_k)
 +\frac1\eta\KL(\pi_k\Vert\pi_{k+1}),\\
 \sum_{k=0}^{\infty}\KL(\pi_k\Vert\pi_{k+1})
 &\leq\eta\KL(\pi_0\Vert\pi_T).
\end{align*}
后一式由前式对有限轮数求和、去掉非负项，再令轮数趋于无穷得到。它控制整个 hard-sync 轨迹的累计 KL 移动量。

Gibbs target 对任意正步长仍是一个严格凸目标的唯一最优点。具体地，定义
\[
 F_\eta(p)=\KL(p\Vert\pi_T)
 +(\eta^{-1}-1)\KL(p\Vert\pi).
\]
直接代入 $\pi^+$ 的 log-density，得到
\[
 F_\eta(p)=\eta^{-1}\KL(p\Vert\pi^+)
 -\eta^{-1}\log Z_\eta(\pi).
\]
所以 $\eta>1$ 时虽然第二个 KL 的系数为负，整个目标的熵系数仍为 $1/\eta>0$，最优点仍唯一。最优点存在、hard-sync 稳定和 teacher KL 单步下降是三个不同结论。

\subsection{KL 收缩的两个数值反例}

取二元词表，$\pi_0=(10^{-12},1-10^{-12})$、$\pi_T=(0.1,0.9)$、$\eta=1.1$。直接代入更新得到
\[
 \pi_1=(0.5856848855168,\;0.4143151144832),
\]
\[
 \KL(\pi_0\Vert\pi_T)=0.1053605156314,
 \qquad\KL(\pi_1\Vert\pi_T)=0.7138510884672.
\]
跨度严格缩小到初值的 $0.1$，正向 KL 却增加到约 $6.7753$ 倍。这说明 $1<\eta<2$ 的稳定性不能表述为每一步 teacher KL 下降。

取 $\pi_0=(1-10^{-12},10^{-12})$、$\pi_T=(0.01,0.99)$、$\eta=0.5$，得到
\[
 \pi_1=(0.9999900502246,\;0.0000099497753721),
\]
\[
 \KL(\pi_0\Vert\pi_T)=4.6051701859549,
 \qquad\KL(\pi_1\Vert\pi_T)=4.6049999147314.
\]
此时 KL 的比值约为 $0.9999630$，超过 $1-\eta=0.5$ 与 $(1-\eta)^2=0.25$。这一次 KL 确实下降，但远处的下降速度不能套用 teacher 附近的二阶渐近速率。

\subsection{从跨度到分布误差的最优 TV 界}

\begin{theorem}[Hilbert 距离控制 TV 的 sharp 界]
\label{thm:hs-sharp-tv}
任意严格正概率分布 $p,q$ 满足
\begin{equation}\label{eq:hs-sharp-tv}
 \TV(p,q)\leq\tanh\!\left(\frac{d_H(p,q)}4\right).
\end{equation}
对每个给定正跨度，都存在二元分布达到等号。
\end{theorem}

\begin{proof}
令 $L_v=p(v)/q(v)$、$m=\min L_v$、$M=\max L_v$、$\delta=d_H(p,q)$。归一化给出 $\E_qL=1$，所以 $m\leq1\leq M$，且 $M/m=e^\delta$。$\delta=0$ 时两个分布相等，结论成立。以下设 $\delta>0$。

凸函数 $|L-1|$ 在 $[m,M]$ 上不超过端点连接的弦线，因此
\[
 |L-1|\leq\frac{M-L}{M-m}(1-m)
 +\frac{L-m}{M-m}(M-1).
\]
在 $q$ 下取期望，再除以二，得到
\[
 \TV(p,q)=\tfrac12\E_q|L-1|
 \leq\frac{(M-1)(1-m)}{M-m}.
\]
记 $h=e^\delta$，代入 $M=hm$ 将右侧化为
\[
 \frac{h+1-hm-m^{-1}}{h-1}.
\]
由 $hm+m^{-1}\geq2\sqrt h$，继续得到
\[
 \TV(p,q)\leq\frac{\sqrt h-1}{\sqrt h+1}
 =\tanh(\delta/4).
\]
为了达到等号，取两个 token 的概率比分别为 $M=e^{\delta/2}$ 和 $m=e^{-\delta/2}$，在 $q$ 下赋予质量
\[
 q(L=M)=\frac1{1+e^{\delta/2}},\qquad
 q(L=m)=\frac{e^{\delta/2}}{1+e^{\delta/2}}.
\]
这些质量满足 $\E_qL=1$。随机变量只取端点值，使弦线界取等号；$hm=m^{-1}$ 又使算术--几何均值界取等号。故整个界达到等号。
\end{proof}

结合 hard-sync 闭式立即得到
\[
 \TV(\pi_k,\pi_T)\leq\tanh(|A_k|b_0/4).
\]
对于常步长 $0<\eta<2$，若要求 TV 不超过 $\delta\in(0,1)$，充分条件是
\[
 |1-\eta|^k b_0\leq4\operatorname{artanh}(\delta).
\]
当 $b_0$ 已不超过右侧时，初始策略就满足要求；否则取 $\eta\neq1$，可解出所需 outer 轮数
\[
 k\geq
 \frac{\log\!\left(b_0/[4\operatorname{artanh}(\delta)]\right)}
 {-\log|1-\eta|}.
\]
这个保证以逐 prefix 的精确 Gibbs 更新为前提。它没有把有限 inner、参数共享或 rollout 测度变化隐藏到精确公式中。
